\chapter{نخستین برنامه}%
\label{ch:first-program}

فکر کردن را شروع کرده‌ایم؛ حالا وقتِ نوشتن است. در این فصل نخستین برنامهٔ
خود را می‌نویسیم، اجرایش می‌کنیم و با ساختارِ کلیِ هر برنامهٔ سلام آشنا
می‌شویم.

\section{چه چیزهایی لازم داریم؟}

برای برنامه نوشتن فقط دو ابزار لازم است:

\begin{itemize}
  \item \textbf{یک ویرایشگرِ متن:} برنامه‌ها فایل‌های متنیِ ساده‌اند. هر
        ویرایشگری که متنِ ساده ذخیره کند کافی است؛ ویرایشگرهایی مانندِ
        \lr{VS Code} کار را خوشایندتر هم می‌کنند (رنگ‌آمیزیِ کد،
        شماره‌گذاریِ خط‌ها).
  \item \textbf{کامپایلرِ سلام:} همان مترجمی که در فصلِ
        \ref{ch:what-is-programming} گفتیم؛ برنامه‌ای به نامِ \code{salam}.
        فرض می‌کنیم روی رایانهٔ شما آماده است و در خطِ فرمان با نوشتنِ
        \code{salam} در دسترس است. برای آزمایش، این فرمان را در ترمینال
        اجرا کنید تا شمارهٔ نسخه را ببینید:
\end{itemize}

\begin{salamen}[language={}]
salam -v
\end{salamen}

\begin{note}
«ترمینال» یا «خطِ فرمان» همان پنجرهٔ سیاه‌رنگی است که در آن فرمان‌ها را
تایپ می‌کنید و با کلیدِ \lr{Enter} اجرا می‌شوند. در ویندوز برنامهٔ
\lr{Terminal} یا \lr{PowerShell} را باز کنید؛ در لینوکس و مک برنامهٔ
\lr{Terminal} را.
\end{note}

\section{سلام، دنیا!}

سنتی دیرینه میانِ برنامه‌نویسان هست: نخستین برنامه در هر زبانِ تازه،
برنامه‌ای است که عبارتِ «سلام، دنیا!» را نمایش می‌دهد. ما هم همین سنت را
پاس می‌داریم. فایلی به نامِ \code{hello.salam} بسازید و این سه خط را در آن
بنویسید:

\begin{salam}
تابع اصلی:
    چاپ "سلام، دنیا!"
تمام
\end{salam}

حالا در ترمینال، در همان پوشه‌ای که فایل را ذخیره کرده‌اید، بنویسید:

\begin{salamen}[language={}]
salam exec hello.salam
\end{salamen}

و خروجی:

\begin{progoutfa}
سلام، دنیا!
\end{progoutfa}

تبریک! شما رسماً برنامه‌نویس شدید. حالا این سه خط را کالبدشکافی کنیم.

\section{کالبدشکافیِ برنامه}

\subsection{تابعِ اصلی: نقطهٔ شروع}

خطِ \code{تابع اصلی:} به سلام می‌گوید: «اجرای برنامه از اینجا شروع می‌شود.»
هر برنامهٔ سلام باید دقیقاً یک \code{تابع اصلی} داشته باشد. واژهٔ
\code{تمام} در خطِ آخر هم پایانِ آن را نشان می‌دهد. فعلاً این دو خط را
مانندِ جلدِ یک دفتر ببینید: هر چه می‌نویسیم، میانِ این دو خط می‌نویسیم. در
فصلِ \ref{ch:functions} خواهیم دید «تابع» چیزهای بیشتری در چنته دارد.

\subsection{تورفتگی: نظمِ دیداریِ برنامه}

دقت کنید که خطِ میانی کمی جلوتر از دو خطِ دیگر نوشته شده است. به این
جلوآمدگی \textbf{تورفتگی} می‌گویند و نشان می‌دهد این خط \emph{درونِ} تابعِ
اصلی است. عادت کنید هر چه درونِ چیزی است را با چهار فاصله تو ببرید؛
برنامه‌تان بسیار خواناتر می‌شود.

\subsection{چاپ: نمایشِ خروجی}

دستورِ \code{چاپ} هر چه به او بدهید را نمایش می‌دهد و به خطِ بعد می‌رود.
متنِ موردِ نظر را میانِ دو علامتِ نقل‌قول (\code{"}) می‌گذاریم؛ به چنین
متنی \textbf{رشته} می‌گویند. می‌توانید چند چیز را با کاما از هم جدا کنید تا
پشتِ سرِ هم چاپ شوند، و فقط عدد و حاصلِ محاسبه هم می‌توانید بدهید:

\begin{salam}
تابع اصلی:
    چاپ "دو به‌علاوهٔ سه می‌شود:", 2 + 3
    چاپ "سال:", 1404
تمام
\end{salam}

\begin{progoutfa}
دو به‌علاوهٔ سه می‌شود: 5
سال: 1404
\end{progoutfa}

به خطِ اول دقت کنید: \code{"دو به‌علاوهٔ سه می‌شود:"} چون درونِ نقل‌قول
است، همان‌طور که هست چاپ شد؛ اما \code{2 + 3} چون بیرونِ نقل‌قول است،
\emph{محاسبه} شد و حاصلش چاپ گشت. این تفاوتِ مهمی است: نقل‌قول یعنی «عینِ
همین متن»، بی‌نقل‌قول یعنی «حسابش کن».

در کنارِ \code{چاپ}، دستورِ \code{بنویس} هم هست که همان کار را می‌کند اما
به خطِ بعد \emph{نمی‌رود}؛ نوشتهٔ بعدی در ادامهٔ همان خط می‌آید:

\begin{salam}
تابع اصلی:
    بنویس "سلام، "
    بنویس "دنیا"
    چاپ "!"
تمام
\end{salam}

\begin{progoutfa}
سلام، دنیا!
\end{progoutfa}

\section{توضیح‌ها: یادداشت برای انسان‌ها}

گاهی می‌خواهیم میانِ کد، برای خودمان یا دیگران یادداشتی بگذاریم. هر چه بعد
از \code{//} بیاید، \textbf{توضیح} است: کامپایلر آن را کاملاً نادیده
می‌گیرد و فقط برای خوانندهٔ انسانی است.

\begin{salam}
// این برنامه مساحت یک اتاق را حساب می‌کند.
تابع اصلی:
    چاپ "مساحت اتاق:", 4 * 3  // طول ۴ متر، عرض ۳ متر
تمام
\end{salam}

\begin{tip}
توضیحِ خوب نمی‌گوید کد \emph{چه} می‌کند (خودِ کد که جلوی چشم است)؛ می‌گوید
\emph{چرا} این کار را می‌کند. مثلاً «طول ۴ متر، عرض ۳ متر» اطلاعی می‌دهد که
از خودِ \code{4 * 3} فهمیده نمی‌شود.
\end{tip}

\section{وقتی اشتباه می‌کنیم}

بیایید عمداً خرابکاری کنیم تا از خطا نترسیم. نقل‌قولِ پایانیِ رشته را جا
بیندازید:

\begin{salam}
تابع اصلی:
    چاپ "سلام، دنیا!
تمام
\end{salam}

کامپایلر برنامه را اجرا نمی‌کند و در عوض پیامی می‌دهد که می‌گوید در کدام
\emph{فایل} و کدام \emph{خط} چه چیزی اشکال دارد (اینجا: رشته‌ای که باز شده
اما بسته نشده است). خطا خواندن مهارتی است که به‌مرور در شما شکل می‌گیرد؛
از همین حالا عادت کنید پیامِ خطا را با حوصله بخوانید و اولین سرنخ را همان
شمارهٔ خط بگیرید.

\begin{warn}
علامت‌های نگارشیِ کد باید همان علامت‌های انگلیسیِ صفحه‌کلید باشند:
نقل‌قولِ \code{"} و کامای \code{,}. اگر صفحه‌کلیدِ فارسیِ شما گیومهٔ
«فارسی» یا ویرگولِ فارسی (،) تایپ کند، کامپایلر آن را نمی‌شناسد و خطا
می‌گیرد. متنِ \emph{درونِ} نقل‌قول‌ها آزاد است و هر نویسه‌ای می‌تواند
باشد.
\end{warn}

\section{ساختنِ فایلِ اجرایی}

فرمانِ \code{salam exec} برنامه را همان لحظه اجرا می‌کند و برای یادگیری و
آزمایش عالی است. راهِ دوم، ساختنِ یک فایلِ اجراییِ مستقل است که بدونِ
کامپایلر هم اجرا می‌شود و سریع‌تر است:

\begin{salamen}[language={}]
salam build hello.salam --output=hello.exe
./hello.exe
\end{salamen}

در این کتاب همیشه از \code{salam exec} استفاده می‌کنیم؛ ساده‌تر است و برای
برنامه‌های آموزشی تفاوتی حس نمی‌کنید.

\section{تمرین‌ها}

\exercise برنامه‌ای بنویسید که نامِ شما را در یک خط و شهرتان را در خطِ بعد
چاپ کند.

\exercise برنامه‌ای بنویسید که حاصلِ \code{7 * 8} را همراهِ توضیحی خوانا
چاپ کند؛ خروجی چنین باشد: \code{حاصل ضرب هفت در هشت: 56}.

\exercise فقط با \code{بنویس} و یک \code{چاپ}، عبارتِ
\code{یک، دو، سه!} را در یک خط چاپ کنید.

\exercise عمداً واژهٔ \code{تمام} را از برنامهٔ سلام دنیا حذف کنید و اجرا
کنید. پیامِ خطا چه می‌گوید؟ به کدام خط اشاره می‌کند؟

\exercise در دستورِ \code{چاپ "2 + 3"} علامت‌های نقل‌قول را برداریم چه
تغییری در خروجی رخ می‌دهد؟ اول حدس بزنید، بعد امتحان کنید.
